% Lab 7 report — two columns, short.
% User-testing results come from the filled forms user-testing-form-P1..P5-filled.pdf.
% Compile: pdflatex test-report.tex

\documentclass[10pt,a4paper,twocolumn]{article}
\usepackage[utf8]{inputenc}
\usepackage[T1]{fontenc}
\usepackage[margin=1.7cm,columnsep=0.7cm]{geometry}
\usepackage{lmodern}
\usepackage{microtype}
\usepackage{array}
\usepackage{booktabs}
\usepackage{tabularx}
\usepackage{enumitem}
\usepackage{amsmath}
\usepackage[table]{xcolor}
\usepackage[hidelinks]{hyperref}
\usepackage{titlesec}

\titlespacing*{\section}{0pt}{8pt}{3pt}
\titlespacing*{\subsection}{0pt}{6pt}{2pt}
\setlist{itemsep=1pt,topsep=2pt,leftmargin=*}
\setlength{\parskip}{3pt}
\setlength{\parindent}{0pt}
\renewcommand{\arraystretch}{1.15}
\raggedbottom

\newcommand{\todo}[1]{\textcolor{red!70!black}{\textbf{[#1]}}}
\newcommand{\blank}{\rule{0.7cm}{0.4pt}}
\newcommand{\fail}{\textcolor{red!70!black}{\textbf{Fail}}}
\newcommand{\partres}{\textcolor{orange!80!black}{\textbf{Part.}}}
\newcommand{\pass}{\textcolor{green!45!black}{\textbf{Pass}}}
\newcommand{\helped}{\textcolor{orange!80!black}{\textbf{Help}}}

\title{\vspace{-1.2cm}\textbf{Lab 7 Report: Usability Evaluation}\\[2pt]
  \large MMAC Study Hub and Aura Screen}
\author{Nongmaithem Hans \quad Himanshu Patel}
\date{\vspace{-0.4cm}}

\begin{document}
\maketitle

%---------------------------------------------------------------------
\section{Introduction}
We evaluated two products from the Multimodal Affective Computing (MMAC) study.
The \textbf{Study Hub} is a website that introduces the study and showcases its
thermal and ECG data in interactive dashboards. \textbf{Aura Screen} is a patient
health-screening app for phones and screening kiosks: sign up, scan at a station,
read the report.

Three problems drove this evaluation. Aura Screen was built for a desktop browser,
so its phone/kiosk view could only be shown through browser developer tools, and
screens needed scrolling. The hub linked to separate Streamlit sensor dashboards
that stopped launching after a while because of deprecated dependencies. And the
project's goal moved from \emph{recording} data to \emph{showcasing} it, which the
recording dashboards did not fit.

%---------------------------------------------------------------------
\section{Test plan}
\textbf{Tasks.}
\begin{itemize}
  \item \textbf{T1 (Aura Screen):} create an account, finish onboarding,
    start a scan and read your heart rate on the report.
  \item \textbf{T2 (Study Hub):} find the ECG data, compare heart rate at rest
    vs city driving, and start a data-access request.
\end{itemize}
T1 is the app's core end-to-end journey; T2 covers the hub's two jobs, showing data
and controlling its release, and requires reading a chart.

\textbf{Users.} Five people from outside the team (P1--P5), tested on 28--29
September 2026. All five were technical/research students, the group closest to the
hub's researcher audience; only P4 had used a similar health or data app before.

\textbf{Procedure.} About 12 minutes each: consent, a think-aloud warm-up, T1, T2,
the SUS questionnaire and two debrief questions. One of us facilitated; the other
observed and filled in the test form. All sessions ran on a laptop; Aura Screen
was used in its phone-sized view in the browser, not on a real phone.

\textbf{Metrics.} Completion with or without help (effectiveness), think-aloud
confusion points (why), and SUS (standard satisfaction score; 68 is average). The
form also had time and error fields, but these were not filled in during the
sessions, so no timing or error counts are reported.

%---------------------------------------------------------------------
\section{Findings}
\subsection{User testing}\label{sec:ut}
\begin{center}\small
\begin{tabular}{@{}lccc@{}}
\toprule
 & T1 Aura Screen & T2 Study Hub & SUS \\
\midrule
P1 & \helped & \helped & 77.5 \\
P2 & \pass & \pass & 62.5 \\
P3 & \helped & \helped & 67.5 \\
P4 & \pass & \pass & 90.0 \\
P5 & \pass & \helped & 67.5 \\
\midrule
Completed & 5/5 & 5/5 & \\
Unaided & 3/5 & 2/5 & \textbf{73.0} \\
\bottomrule
\end{tabular}\\[2pt]
{\footnotesize \pass{} = completed alone; \helped{} = completed with help.}
\end{center}
Every participant finished both tasks; nobody failed. Half of the attempts needed a
hint, most of them on the website (3 of 5). The mean SUS of \textbf{73.0} (range
62.5--90) is above the 68 average. SUS rated both products together. The lowest
item was confidence (Q9, mean 3.6/5), matching the hints needed.

\textbf{Aura Screen (T1).} P2, P4 and P5 called it easy (``self-explanatory'';
P4 liked the step bar). P1's browser blocked the camera: P1 was stuck until shown
the \emph{Continue without camera} button (U1). P3 missed the consent tick box in
the long record form and had to go back (U2). P5 paused to read the scan
instructions.

\textbf{Study Hub (T2).} The resting heart rate was found quickly (P1: 79\,bpm on
the top card). The city-driving rate was not shown as a number, so P1 and P3 had to
read it off the chart (U3). P3 found the ECG page only through the Modalities menu,
and P5 did not see that the menu was clickable (U4). P2 did not know what ECG meant
at first (U5). P1 found scrolling past the large charts slow (U6). On the access
request, P4 was unsure whether the supervisor fields applied to them (U7), and P5
found the agreement text long and dense (U8). P2 did not realise that apps such as
Aura Screen are built on the data (U9). P4 found the three-step form and
draggable widgets clear.

\begin{table*}[t]
\centering\small
\caption{Heuristic evaluation: merged findings from both evaluators (inspection).}
\label{tab:heuristic}
\begin{tabularx}{\textwidth}{@{}l>{\raggedright\arraybackslash}p{3.3cm}X>{\raggedright\arraybackslash}p{4.6cm}c@{}}
\toprule
ID & Screen / step & Problem & Heuristic violated & Sev. \\
\midrule
\multicolumn{5}{@{}l}{\textbf{Aura Screen}} \\
A1 & Any screen, desktop browser & Phone/kiosk view only reachable through browser
  developer tools. & 7. Flexibility \& efficiency; 4. Consistency \& standards &
  \textbf{3} \\
A2 & Scan Hub (T1, start scan) & Core screens scroll in the phone view: Scan Hub is
  +454\,px taller than a 390$\times$844 screen. & 8. Aesthetic \& minimalist design &
  \textbf{3} \\
A3 & Session report (T1, read HR) & Third heart-rate tile cut off at phone width. &
  8. Aesthetic \& minimalist design & \textbf{3} \\
A4 & Navigation, all screens & Navigation sits at the top, out of thumb reach,
  unlike phone-app convention. & 4. Consistency \& standards & 2 \\
A5 & Sign-up, record setup (T1) & 13\,px inputs make iOS zoom in on every tap. &
  4. Consistency \& standards & 2 \\
A7 & Kiosk layout, all screens & Kiosk screens leave large empty areas. &
  8. Aesthetic \& minimalist design & 2 \\
A6 & Buttons and cards (touch) & Hover styles ``stick'' after a tap, so a tapped item
  looks still selected. & 1. Visibility of system status & 1 \\
\midrule
\multicolumn{5}{@{}l}{\textbf{Study Hub website}} \\
W1 & Home $\rightarrow$ sensor dashboard link (T2, find data) & Links lead nowhere
  when a dashboard server is down or fails to launch (deprecated versions); no
  message explains why. & 9. Recover from errors; 1. Visibility of system status &
  \textbf{4} \\
W2 & Sensor dashboards (T2, compare HR) & Dashboards are recording tools; visitors
  only want to view data. & 2. Match with the real world & \textbf{3} \\
W3 & Sensor dashboards (T2) & Separately styled Streamlit apps look and work
  differently from the site. & 4. Consistency \& standards & \textbf{3} \\
W4 & Opening the site & Site fails to start (stray build cache); no run
  instructions. & 10. Help \& documentation & \textbf{3} \\
W5 & Home page sections & Placeholder ``lorem ipsum'' text. &
  8. Aesthetic \& minimalist design & 2 \\
W6 & Data download (T2, request data) & Data offered as open downloads with no
  terms of use. & 5. Error prevention & 2 \\
W7 & ECG charts (T2, compare HR) & Charts work with a mouse only. &
  7. Flexibility \& efficiency & 2 \\
W9 & Request access, submit (T2) & Sending relies on an email app being set up;
  nothing happens without one. & 9. Recover from errors & 2 \\
W8 & ECG dashboard widgets & Widgets can be dragged, but nothing shows this. &
  6. Recognition rather than recall & 1 \\
W10 & Thermal page & Sample temperatures are illustrative, not measured. &
  2. Match with the real world & 1 \\
\bottomrule
\end{tabularx}
\end{table*}

\begin{table*}[t]
\centering\small
\caption{Problems found only in user testing (Section~\ref{sec:ut}), rated on the same scale.}
\label{tab:ut}
\begin{tabularx}{\textwidth}{@{}l>{\raggedright\arraybackslash}p{3.3cm}X>{\raggedright\arraybackslash}p{4.6cm}c@{}}
\toprule
ID & Screen / step & Problem (participants) & Heuristic violated & Sev. \\
\midrule
U1 & Scan step (T1, start scan) & Camera blocked; the way past it (\emph{Continue
  without camera}) is easy to miss, so the scan cannot start (P1). &
  9. Recover from errors; 1. Visibility of system status & \textbf{3} \\
U3 & ECG dashboard (T2, compare HR) & Per-phase heart rate (rest vs city) not shown
  as a number; users read it off the chart (P1, P3). &
  6. Recognition rather than recall & \textbf{3} \\
U2 & Record setup (T1) & Long form; the consent box was skipped and only caught on
  submit (P3). & 5. Error prevention & 2 \\
U4 & Home $\rightarrow$ Modalities menu (T2, find data) & Modality pages hard to
  find; the menu does not look clickable (P3, P5). &
  6. Recognition rather than recall & 2 \\
U5 & ECG dashboard (T2) & Jargon (``ECG'') and charts with no caption saying what
  they show (P2, P3). & 2. Match with the real world; 10. Help \& documentation & 2 \\
U6 & ECG dashboard, scrolling & Scrolling lags past the large charts (P1). &
  7. Flexibility \& efficiency & 2 \\
U7 & Request access, details step (T2) & Unclear which fields (supervisor) apply to
  whom (P4). & 5. Error prevention & 2 \\
U8 & Request access, agreement step (T2) & Agreement is a wall of legal text with no
  summary (P5). & 8. Aesthetic \& minimalist design & 2 \\
U9 & Home page & Apps built on the data are easy to miss; the hub reads as
  ``research only'' (P2). & 8. Aesthetic \& minimalist design & 1 \\
\bottomrule
\end{tabularx}
\end{table*}

\subsection{Heuristic evaluation}
\textbf{Method.} Each of us walked through T1 and T2 on our own and noted every
problem against Nielsen's ten heuristics. For each problem we recorded where it
occurs (screen and task step), the heuristic it violates and a severity: 0 not a
problem, 1 cosmetic, 2 minor, 3 major, 4 catastrophe. We then merged the two lists
and removed duplicates. The inspection covered the builds from before this lab's
changes. Everything was checked on a laptop, so A4--A6, which only affect touch
screens, were judged from the design rather than observed.

\textbf{Results.} The merged list has 17 problems (Table~\ref{tab:heuristic}): one
catastrophe (W1), six major, seven minor and three cosmetic. The most-violated
heuristics were 8 (minimalist design) and 4 (consistency), with four problems each.
Most came from Aura Screen being a desktop layout squeezed into a phone-sized
screen. No violations of 3
(user control and freedom) were found.

\textbf{Inspection vs user testing.} The two methods found different problems.
Inspection caught layout, platform and robustness problems that users seldom
mention, such as dead links, iOS zoom and sticking hover styles. Testing found nine
further problems (Table~\ref{tab:ut}). These were mostly about understanding: the
missing per-phase numbers, jargon, the dense agreement and the blocked camera. Only
one problem was close to an inspection finding: charts being hard to use (W7) and
reading values off a chart (U3).

\subsection{Accessibility check}
A quick manual check of both products, looking for what a user would notice. We
read each screen on a laptop (Aura Screen in its phone-sized view), used the
website with the keyboard only,
and checked that images and icon buttons are labelled. No tools or measurements.

\begin{center}\small
\begin{tabularx}{\columnwidth}{@{}lXc@{}}
\toprule
Check & What we saw & \\
\midrule
Text, website & Main text is clear, but faded grey text (inactive menu
  items, notes on cards) is hard to read. & \fail \\
Text, charts & Axis labels are small and pale; hard to read at a glance. &
  \fail \\
Text, Aura & Body text and buttons are clear; small captions and pink accent
  text are faint. & \partres \\
Text size & Some labels and badges are very small, especially on charts. & \partres \\
Keyboard & Tab reaches every widget and Alt+arrows moves it; the request dialog
  keeps focus and Esc closes it; charts cannot be explored without a mouse. &
  \partres \\
Labels & Images have descriptions, but the theme toggle is an unlabelled icon, so a
  screen reader cannot name it. & \partres \\
\bottomrule
\end{tabularx}
\end{center}

\subsection{Live demo}
Aura Screen on a laptop: the old build needs dev-tools emulation and Scan Hub scrolls;
the phone/kiosk frame fits the same screen without scrolling (A1, A2).

%---------------------------------------------------------------------
\section{Improvement plan}
Ranked by severity, then by how many participants hit the issue. ``Ready'' =
prototype built and used in these sessions; ``Planned'' = next iteration.

\textbf{Ready:} W1--W3 showcase pages replace the Streamlit links; A1--A5, A7
phone/kiosk-first Aura Screen (device frame, no-scroll screens, bottom tabs,
16\,px inputs, fixed report); W4 one-command launcher; W6 access request with a
data-use agreement.

\begin{center}\footnotesize
\begin{tabularx}{\columnwidth}{@{}p{1.25cm}X@{}}
\toprule
Issues & Planned change (in priority order) \\
\midrule
U3 & Print mean heart rate per phase (rest, city) beside the chart. \\
U1 & Explain a blocked camera up front; \emph{Continue without camera} as the main
  button. \\
Readability & Darken faded grey, caption and chart-label text in both products. \\
U4, U5 & Modality cards on the home page; a one-line caption under each chart; a
  plain-language line on what ECG/thermal measure. \\
U7, U8 & Mark optional fields; supervisor fields for students only; bullet summary
  above the agreement. \\
U2 & Flag missed required fields (consent) in place. \\
U6 & Render off-screen charts lazily. \\
W7, kbd & Keyboard access and a data table per chart. \\
U9, labels & Link ``Built with the data'' from the hero; name the theme toggle;
  enlarge the smallest text. \\
\bottomrule
\end{tabularx}
\end{center}
\textbf{Limitations.} All five participants were technical students, so
non-technical (patient-like) users still need testing. Everything ran on a laptop,
so Aura Screen still needs testing on real phones and kiosks. Task times and errors
were not recorded. We will address all three in the next round, when we re-test the
Planned changes.

\end{document}
